\section*{ID9255: Definition of dˆ3r’Agyro}
\paragraph{Metadata.}
PiID: 6821207388377 | RTEID: RTE:P35939.2203:T2.6681 | UUID: 236f942f-bd82-52a5-9885-da4eea3fc0d8

\paragraph{Canonical Statement.}
ributions: Atotal(r,t)=Adipole+Agyro+AinducedA\_\{\textbackslash{}text\{total\}\}(r,t) = A\_\{\textbackslash{}text\{dipole\}\} + A\_\{\textbackslash{}text\{gyro\}\} + A\_\{\textbackslash{}text\{induced\}\}Atotal(r,t)=Adipole+Agyro+Ainduced, where \$Adipole=(µ0/4π)µeff(t)×\$ (rrcm)/rrcm3A\_\{\textbackslash{}text\{dipole\}\} = (\textbackslash{}mu\_0/4\textbackslash{}pi) \textbackslash{}mu\_\{\textbackslash{}text\{eff\}\}(t) \textbackslash{}times (r - \$r\_\{\textbackslash{}text\{cm\}\})/|r-r\_\{\textbackslash{}text\{cm\}\}|ˆ3Adipole=(µ0/4π)µeff(t)×\$ (rrcm)/rrcm3, \$Agyro=(µ0/4π)\$ \$∫\$ VgyroJgyro(r,t)× (rr)/rr3d3rA\_\{\textbackslash{}text\{gyro\}\} = (\textbackslash{}mu\_0/4\textbackslash{}pi) \textbackslash{}int\_\{V\_\{\textbackslash{}text\{gyro\}\}\} J\_\{\textbackslash{}text\{gyro\}\}(r’,t) \textbackslash{}times (r - r’)/|r - r’|ˆ3 ; \$dˆ3r’Agyro=(µ0/4π) ∫\$ VgyroJgyro(r,t)× (rr)/rr3d3r, and \$Ainduced=(1/4πc2)\$ \$∫ V(∂J/∂t)×\$ (rr)/rrd3rA\_\{\textbackslash{}text\{induced\}\} = -(1/4\textbackslash{}pi cˆ2) \textbackslash{}int\_V (\textbackslash{}partial J/\textbackslash{}partial t) \textbackslash{}times (r - r’)/|r - r’| ; \$dˆ3r’Ainduced=(1/4πc2) ∫\$ \$V(∂J/∂t)×\$ (rr)/rrd3r. Interpretive Meaning: The vector potential is decomposed into contributions from the effective dipole, the gyroscopic current distribution, and induced currents due to timevar

\paragraph{Canonical Equation.}
\[
dˆ3r’Agyro=(µ0/4π) ∫
\]

\paragraph{Dictionary.}
Definition of dˆ3r’Agyro — dˆ3r’Agyro=(µ0/4π) ∫

\paragraph{Encyclopedia.}
Definition of dˆ3r’Agyro is an algebraic definition, written dˆ3r’Agyro=(µ0/4π) ∫. It is an algebraic definition expressing the quantity directly in terms of the variables on its right-hand side. Within the Trawin Zero Point registry it serves as one element of the framework's mathematical narrative.
